% This is an includable section; the parent report supplies theorem environments.
\section{A source invariant permits additional exact zero tests}
\label{sec:runtime-known-zero}

The prior product certificate associates with each instruction leaf $i$ nine
nonempty masks $m_{i,r}\subseteq\{0,1,2\}$, or marks the leaf unreachable.
The abstract labels $0,1,2$ denote the numerical sets $\{0\}$, $\{1\}$, and
$\{2,3,\ldots\}$, respectively.  Its initial condition is $T=W=H=0$, with the other six
logical registers arbitrary.  The checker reparses the original program and
checks abstract transfer across every feasible edge; equality of all $8391$
labelled control-flow edges transports this certificate to each new layout.
The two verified blank bootstraps have logical registers
$(3,0,0,0,3,0,0,0,0)$ and therefore satisfy this initial condition.

The affine box denotes the entire encoded register file, including its finite
nonempty spare suffix.  In the inherited interpreter, a box
$\operatorname{plain}(\varnothing)$ denotes the unchanged input file.  Marking
register $r$ is permitted only in this state and produces
$\operatorname{mark}(r)$.  An increment or positive decrement produces
$\operatorname{plain}(e)$ for a nonempty effect $e$; this state cannot be marked
again.  The existing zero-decrement contract, starting in state $95$, restores
the entire file and ends in state $112$, with no extra left-boundary cell.

\paragraph{Additional rule.}
At the start of a guarded instruction, intersect the selected register's mask
with the actual zero or positive decrement guard.  Reject an empty
intersection.  Convert the resulting masks from logical to physical register
order using $(6,2,0,4,3,7,8,5,1)$.  In addition to all inherited oracle rules,
allow
\[
 (95,\operatorname{mark}(r))\longmapsto
 (112,\operatorname{plain}(\varnothing),0)
 \qquad\text{if }m_r=\{0\}.
\]
The final zero is the additional left-boundary cell count.  This rule applies
even when $r$ differs from the register named by the source instruction.

\paragraph{Soundness.}
Induct on the number of oracle calls before the unique nonempty effect.
Initially, the boxed file is exactly the source input.  A mark operation only
marks that file.  If the new rule applies, the selected value is zero by the
source mask and the induction hypothesis.  The inherited zero-decrement
contract therefore restores precisely the same entire file, including all
other registers and all spares, and returns to the unchanged box state.
The induction hypothesis is preserved.  After any nonempty increment or
positive-decrement effect, marking is disabled, so this rule cannot be
misapplied using a stale source mask.  Literal movement over the program-counter
region and the boxed boundary uses the unchanged affine rules.  Consequently a
successful exact-target derivation remains valid for every gap parameter and
every complete file satisfying the product masks and the selected guard.

The $42$ literal transition rows supporting the register and marking induction
contracts must remain byte-for-byte equal to the Context246 baseline, with
their targets intact.  The implementation enforces this condition before
checking a candidate.  The endpoint still specifies the complete register file,
all spares, the head, the program-counter residue, and the gap change; the rule
does not permit unspecified garbage or ignore additional side effects.

\paragraph{Limits of the result.}
This is a compositional mathematical extension of the existing affine rule
system, implemented as an ordinary Python checker.  It has not been integrated
into a proof-assistant kernel in this archive.  The four concrete vectors per
guard used in a separate finite merger screen are independent rejection
fixtures; they do not prove this all-parameter rule.  Neither this rule nor
those screens supply a machine with fewer than $246$ states.  Source-mask
admissibility also does not imply that a particular register vector is reachable
from the actual blank bootstrap, because the product invariant overapproximates
reachability and forgets correlations.
